\section{Daubert Standard Analysis}
\label{sec:daubert}

\subsection{The Daubert Framework}

In \textit{Daubert v.\ Merrell Dow Pharmaceuticals}~\citep{daubert1993},
the Supreme Court established four criteria for admissibility of
scientific expert testimony in federal courts:
\begin{enumerate}
\item \textbf{Testability}: The theory or technique can be and has been tested.
\item \textbf{Peer review and publication}: The theory or technique has been
  subjected to peer review and publication.
\item \textbf{Known or potential error rate}: The known or potential error
  rate of the technique is acceptable.
\item \textbf{General acceptance}: The theory or technique is generally
  accepted in the relevant scientific community.
\end{enumerate}

\textit{Daubert} applies directly to federal proceedings; most states
have adopted similar standards.

\subsection{The Audit Chain Under Daubert}

\paragraph{Criterion 1: Testability.}
The \bisect{} audit chain is maximally testable: any party with access to
the public manifest and an internet connection can run \texttt{bisect label-verify}
and independently confirm or refute the original expert's computation.
This exceeds the typical testability requirement: most scientific techniques
require specialized equipment or training to test, whereas \texttt{bisect label-verify}
can be run by any qualified expert in redistricting or computer science.
\textbf{Criterion 1: Satisfied.}

\paragraph{Criterion 2: Peer review and publication.}
The SHA-256 algorithm is standardized by NIST in FIPS PUB 180-4~\citep{fips180}
and has been subjected to extensive peer review and cryptanalysis over two
decades.
The \bisect{} algorithm is documented in B.1~\citep{b1paper} and subsequent
papers, and the audit chain architecture is described in this paper.
\textbf{Criterion 2: Satisfied.}

\paragraph{Criterion 3: Known error rate.}
The false non-detection rate of the SHA-256 audit chain is bounded by
$2^{-80}$ (probability that a tampered manifest passes verification under
the birthday attack model).
This error rate is $10^{-24}$, far below any reasonable threshold for
legal evidence (which typically requires only that evidence be ``more
probable than not'' or, for criminal cases, ``beyond reasonable doubt'').
\textbf{Criterion 3: Satisfied.}

\paragraph{Criterion 4: General acceptance.}
SHA-256 is the most widely deployed cryptographic hash function in the world,
used in TLS/SSL, code signing, version control (git), and government
document integrity verification.
It is accepted by NIST, FIPS, NSA, and the cryptographic research community.
\textbf{Criterion 4: Satisfied.}

\subsection{The Reproducibility Advantage}

Beyond satisfying Daubert's four criteria, the audit chain provides a
reproducibility guarantee that is unusual in redistricting expert testimony.
Typically, an expert's computational analysis is reproducible in principle
but impractical in practice: the opposing party would need to obtain the
same software, same data, and same parameter settings, and then run the same
analysis.

With \bisect{}'s audit chain, reproducibility is:
\begin{itemize}
\item \textbf{Automatic}: \texttt{bisect label-verify} performs the reproduction.
\item \textbf{Fast}: Under 60 seconds for any state (see Section~\ref{sec:demo}).
\item \textbf{Exact}: The reproduction is bit-for-bit identical; there is no
  ``reproduction with minor variations.''
\item \textbf{Open}: The \bisect{} binary is open-source; the algorithm and
  its implementation can be audited by any qualified expert.
\end{itemize}

This makes the expert's computation ``reproducibly falsifiable''---the
opponent does not need to dispute the expert's methodology; they need only
run \texttt{bisect label-verify} and check whether it passes or fails.
